\subsection{Contexto: O Radiotelescópio BINGO}

\subsubsection{O que é o BINGO?}

O \textbf{BINGO} (Baryon Acoustic Oscillations in Neutral Gas Observations) é um radiotelescópio de última geração projetado para realizar a primeira detecção de Oscilações Acústicas de Bárions (BAO) na faixa de radiofrequência.

\paragraph{Características Técnicas}

\bigskip\noindent
\begin{tabular}{p{\dimexpr 0.500\linewidth-2\tabcolsep}p{\dimexpr 0.500\linewidth-2\tabcolsep}}
\toprule
Característica & Especificação \\
\hline
\textbf{Tipo} & Radiotelescópio de mapeamento de intensidade \\
\textbf{Frequência} & 980 - 1260 MHz \\
\textbf{Espelho primário} & 40 metros de diâmetro \\
\textbf{Espelho secundário} & 35 metros de diâmetro \\
\textbf{Receptores} & 28 cornetas no plano focal \\
\textbf{Campo de visada} & 15° em declinação \\
\textbf{Operação} & Drift-scan (sem partes móveis) \\
\bottomrule
\end{tabular}

\bigskip

\paragraph{Localização}

O telescópio está sendo instalado na \textbf{Serra do Urubu}, município de \textbf{Aguiar}, no Vale do Piancó, Sertão da Paraíba. A escolha do local considerou:

\begin{itemize}
\item Baixa poluição eletromagnética
\item Ausência de sinal de celular nas proximidades
\item Sem rotas aéreas diretas
\item Colaboração da comunidade local
\item Acesso viável para manutenção
\end{itemize}

\paragraph{Missão Científica}

O BINGO foi concebido para responder algumas das perguntas mais fundamentais da cosmologia contemporânea.

O BINGO investigará:

\begin{enumerate}
\item \textbf{Energia Escura}: Componente misteriosa que acelera a expansão do universo
\item \textbf{Matéria Escura}: Matéria invisível que influencia a estrutura cósmica
\item \textbf{Oscilações Acústicas de Bárions (BAO)}: ``Marcadores padrão'' do universo
\item \textbf{Rajadas Rápidas de Rádio (FRB)}: Eventos energéticos misteriosos
\end{enumerate}

\subsubsection{Colaboração Internacional}

O projeto é uma parceria entre:

\bigskip\noindent
\begin{tabular}{p{\dimexpr 0.333\linewidth-2\tabcolsep}p{\dimexpr 0.333\linewidth-2\tabcolsep}p{\dimexpr 0.333\linewidth-2\tabcolsep}}
\toprule
País & Instituições Principais & Contribuição \\
\hline
\textbf{Brasil} & USP, INPE, UFCG, UFPB & Coordenação científica, instalação, operação \\
\textbf{Reino Unido} & Universities of Manchester, Oxford & Instrumentação, análise de dados \\
\textbf{China} & CAS, Yangzhou University & Espelhos, estrutura mecânica \\
\textbf{Suíça} & FHNW & Eletrônica de receptor \\
\textbf{África do Sul} & SARAO & Troca de expertise em radioastronomia \\
\bottomrule
\end{tabular}

\bigskip

\subsubsection{Significado para o Brasil}

O BINGO representa:

\begin{itemize}
\item \textbf{Primeiro radiotelescópio} dessa magnitude construído no Brasil
\item \textbf{Inovação industrial} com fabricação nacional de componentes
\item \textbf{Formação de recursos humanos} em radioastronomia
\item \textbf{Desenvolvimento regional} no Sertão paraibano
\item \textbf{Inserção internacional} da ciência brasileira em cosmologia
\end{itemize}

\subsubsection{Linha do Tempo}

\bigskip\noindent
\begin{tabular}{p{\dimexpr 0.500\linewidth-2\tabcolsep}p{\dimexpr 0.500\linewidth-2\tabcolsep}}
\toprule
Ano & Marco \\
\hline
2017 & Aprovação da colaboração BINGO \\
2019 & Início das obras no sítio em Aguiar \\
2022 & Inauguração formal do projeto \\
2024 & Chegada dos componentes da China \\
2025-2026 & Montagem e testes do telescópio \\
2027 & Primeiras observações científicas \\
\bottomrule
\end{tabular}

\bigskip

\subsubsection{Recursos}

\begin{itemize}
\item \textbf{Site oficial:} \href{https://bingotelescope.org}{https://bingotelescope.org}
\item \textbf{Coordenação Brasil:} Prof. Elcio Abdalla (USP)
\item \textbf{Coordenação UFCG:} Prof. Edmar Candeia Gurjão (DEE/UFCG)
\item \textbf{Coordenação Hermathena:} Luciano Barosi
\end{itemize}

\subsubsection{É neste contexto que surge o projeto Hermathena, concebido para transformar a presença do radiotelescópio BINGO no sertão paraibano em uma oportunidade ampliada de educação científica, comunicação pública da ciência e desenvolvimento regional.}

\textit{Para mais informações técnicas, consulte a documentação científica da colaboração BINGO.}